\section{Introduction}
Networked control systems (NCS) are integral components of sophisticated robotics and avionics, orchestrating the collaborative function of numerous sensors and actuators to achieve objectives such as locomotion, stabilization, sensor fusion, and position control \cite{gupta2009networked,moyne2007emergence}. Additionally, NCS play a pivotal role in manufacturing settings, interlinking multiple machines for synchronized material handling and production planning \cite{vitturi2019industrial}. This field distinguishes itself from conventional networking paradigms in several key aspects:

\begin{enumerate}
    \item \textbf{Prioritization of Message Efficiency Over Throughput:} Unlike traditional networks where throughput is paramount, NCS prioritize efficient message sizes, typically ranging from three to fifty bytes \cite{gupta2009networked}. The emphasis is on the rapidity of message delivery, with many communications comprising single-packet messages.
    \item \textbf{Critical Need for Deterministic Message Delivery:} In NCS, the stability of control loops is contingent upon the guarantee of message delivery within specific time frames \cite{moyne2007emergence}. This necessitates a network where the round trip times (RTTs) have minimal variability, ensuring reliable and predictable communication.
    \item \textbf{Essential Robustness and Simplified Management:} Robustness in NCS is non-negotiable, requiring the absence of single points of failure \cite{vitturi2019industrial}. Furthermore, these systems are designed to operate independently of network controllers for coordinating port-forwarding configurations. Avoiding central management for route adjustments on all nodes is crucial, as it significantly reduces complexities and variations in RTTs.
\end{enumerate}

In summary, NCS demand a unique networking approach focused on small, swift message exchanges, determinism in delivery times, and a robust, decentralized architecture to maintain system stability and efficiency.

Recent developments in the field of NCS have seen a gradual shift from traditional fieldbuses to more contemporary networking solutions. Currently, the state-of-the-art in NCS predominantly involves the use of Switched Ethernet or its derivatives in proprietary protocols \cite{gupta2009networked,moyne2007emergence}. This transition has been driven by the greater accessibility of Switched Ethernet hardware and its status as an open standard. Such features facilitate the integration of control products from various vendors, enabling customers to build more scalable and complex systems without the need for extensive reconfiguration or adherence to FieldBus specifications.

Despite its widespread adoption, Switched Ethernet falls short in addressing the specialized requirements of NCS. Its design, primarily tailored for general-purpose networking, does not align perfectly with the unique needs of NCS, particularly in areas such as real-time data transmission and network robustness. In Section~\ref{sec:related}, we discuss a key limitation of Switched Ethernet: its inability to support more than one route to any given endpoint in a network graph. This constraint leads to bottlenecks at specific switches and introduces Single Points of Failure (SPoF) in network architectures. Consequently, a failure in a single link can render many endpoints irreversibly inaccessible.

To enhance Message Delivery Determinism, a key performance indicator for NCS, we propose the implementation of a multipath routing strategy. This approach allows for the operation of NCS within highly interconnected graph structures, as opposed to limited spanning-tree configurations. It enables the rerouting of messages around congested nodes instead of accumulating them in extensive buffers.

However, current multipath routing methods, predominantly designed for datacenter environments, rely heavily on network controllers (NCs) to manage and update port-forwarding tables. These NCs hold comprehensive information about the entire network's state. In scenarios like a link failure or a sudden surge in network traffic (for instance, due to a router becoming overloaded), these controllers must reconverge to a new suitable routing topology and disseminate updated routing instructions across the network. This reconvergence takes as little as 50~ms with modern Fast Reroute pre-computation~\cite{cisco_lfa} and up to several hundred milliseconds with vanilla link-state protocols~\cite{eramo2008multi,farkas2009performance}. Such delays, albeit brief, can disrupt the entire network operation temporarily, leading to potential failures in the control loops of an NCS. Therefore, these existing multipath routing methods, given their reliance on network controller-based reconvergence, are not ideally suited for application in NCS.

Taken together, the limitations identified above---single points of failure, spanning-tree bottlenecks, and controller-dependent reconvergence delays---motivate the design of a stateless, multipath routing protocol purpose-built for NCS. To the authors' knowledge, TinyNet is the first protocol to combine the simplicity of Switched Ethernet with the multipath routing benefits found in datacenter networks, without requiring a central controller. It enables multipath capabilities without relying on a network controller for managing forwarding tables, achieved through a port-forwarding algorithm that utilizes each router's \textit{local awareness} of neighboring nodes' packet queues and historical data on packet transmissions observed at the router's ports. We describe the protocol rules and evaluate performance through hardware experiments and simulation.